% Part: turing-machines
% Chapter: undecidability
% Section: trakhtenbrot

\documentclass[../../../include/open-logic-section]{subfiles}

\begin{document}

% SEGMENT OLP-0273-S01

\olfileid{tur}{und}{tra}
\olsection{ட்ராக்டென்ப்ரோட்டின் தேற்றம்}

\begin{explain}
  \olref[rep]{sec} இல் ஒரு டியூரிங் பொறி~$M$ மற்றும் உள்ளீட்டுச்
  சரம்~$w$ காக !!{sentence}s $!T(M,w)$,~$!E(M,w)$ ஆகியவற்றை
  வரையறுத்தோம். பின்னர்,~$M$ உள்ளீடு~$w$ இல் தொடங்கப்பட்டு இறுதியில்
  நின்றால், எனில் மட்டுமே $!T(M,w) \lif !E(M,w)$ ஏற்புடையது என்று
  \olref[ver]{lem:valid-if-halt}, \olref[ver]{lem:halt-if-valid}
  ஆகியவற்றில் காட்டினோம். நிற்றல் சிக்கல் தீர்மானிக்கவியலாதது
  என்பதால், முதல்தர தருக்கத்தின் !!{sentence}s இன் ஏற்புடைமையும்
  நிறைவுறத்தக்கதன்மையும் தீர்மானிக்கவியலாதவை என்பது இதனால்
  பின்வருகிறது
  (\cref{tur:und:uns:thm:decision-prob,tur:und:uns:cor:undecidable-sat}).

  ஆனால் வாக்கியங்களின் ஏற்புடைமையும் நிறைவுறத்தக்கதன்மையும்,
  முடிவுறுவதாகவோ முடிவுறாததாகவோ உள்ள தன்னிச்சையான
  !!{structure}s காக வரையறுக்கப்படுகின்றன. முடிவுறு
  !!{structure} ஒன்றில் ஓர் !!a{sentence} நிறைவுறத்தக்கதா என்பதை
  (அல்லது எல்லா முடிவுறு !!{structure}s இலும் ஏற்புடையதா என்பதைத்)
  தீர்மானிப்பது மேலும் எளிதாக இருக்கலாம் என்று நீங்கள் ஐயுறலாம்.
  அவ்வாறு இல்லை என்று காட்டுமாறு தீர்மானச் சிக்கலின்
  தீர்க்கவியலாமைக்கான நிரூபிப்பைத் தகவமைக்கலாம்.

  முதலில், \olref[ver]{lem:halt-if-valid} இன் நிரூபிப்புக்குத்
  திரும்பிச் சென்றால், உள்ளீடு~$w$ இல் தொடங்கப்படும்போது பொறி~$M$
  செய்வதைச் சரியாக விவரிக்கும்~$!T(M,w)$ இன் மாதிரி~$\Struct M$
  ஒன்றை அங்கு உருவாக்கினோம் என்பதைக் காண்பீர். அம்மாதிரியின் சார்பகம்
  ~$\Nat$, அதாவது முடிவுறாதது. ஆனால்~$M$ உண்மையில் உள்ளீடு~$w$ இல்
  நின்றால், அதே வழியில் முடிவுறு மாதிரி~$\Struct M'$ ஒன்றைக் கட்டலாம்.
  உள்ளீடு~$w$ இல் தொடங்கப்பட்ட~$M$,~$k$ படிகளுக்குப் பிறகு நிற்கிறது
  என்று கொள்க. $n$ என்பது~$k$,~$w$ இன் நீளம் ஆகியவற்றில் பெரியது
  எனக்கொண்டு, கணம் $\{0, \dots, n\}$ ஐச் சார்பகம்~$\Domain{M'}$ ஆக
  எடுக்க; மேலும்
  \[
    \Assign{\prime}{M'}(x) = 
    \begin{cases}
      x + 1 &\text{$x < n$ எனில்}\\
      n &\text{இல்லையெனில்,}
    \end{cases}
  \]
  எனவும், $x < y$ அல்லது $x = y = n$ எனில் மட்டுமே
  $\tuple{x,y} \in \Assign{<}{M'}$ எனவும் வைக்க. மற்றபடி
  $\Struct{M'}$ ஆனது~$\Struct{M}$ போலவே வரையறுக்கப்படுகிறது.
  \olref[ver]{lem:halt-if-valid} இன் நிரூபிப்பில் உள்ளதைப்போலவே,
  ~ $\Struct{M'}$ இன் வரையறையின்படி
  $\Sat{M'}{!T(M,w)}$. மேலும்~$M$ உள்ளீடு~$w$ இல் நிற்கிறது என்று
  கருதியதால், $\Sat{M'}{!E(M,w)}$. ஆகவே~$\Struct{M'}$ என்பது
  $!T(M,w) \land !E(M,w)$ இன் ஒரு முடிவுறு மாதிரி
  ($\lif$ ஐ~$\land$ ஆல் மாற்றியுள்ளோம் என்பதைக் கவனிக்க).
  
  ஒரு நிரூபிப்பின் பாதியை அடைந்துவிட்டோம்:~$M$ உள்ளீடு~$w$ இல்
  நின்றால், $!T(M,w) \land !E(M,w)$ கு ஒரு முடிவுறு மாதிரி உள்ளது
  என்று காட்டியுள்ளோம். வருந்தத்தக்க வகையில், இதன் மறுதலை
  உண்மையல்ல; அதாவது, சில உள்ளீடு~$w$ இல் நிற்காத டியூரிங் பொறிகள்
  இருந்தாலும், $!T(M,w) \land !E(M,w)$ கு ஒரு முடிவுறு மாதிரி இன்னும்
  உள்ளது. எடுத்துக்காட்டாக, ஒற்றை நிலை $q_0$ ஐயும்
  $\delta(q_0,\TMblank) = \tuple{q_0,\TMblank,\TMstay}$ என்ற
  கட்டளையையும் கொண்ட பொறி~$M$ ஐக் கருதுக. வெற்று உள்ளீடு
  ~$w = \emptyseq$ இல் தொடங்கப்பட்டால், இப்பொறி ஒருபோதும் நிற்பதில்லை:
  அது ஒரு முடிவிலிக் கண்ணியில் உள்ளது; ஆனால் நாடாவை மாற்றுவதோ தலையை
  நகர்த்துவதோ இல்லை. எல்லா அமைவுநிலைகளும் ஒரே மாதிரியானவை (அதே நிலை,
  அதே தலை இருப்பிடம், அதே நாடா உள்ளடக்கம்).
  $!T(M,\emptyseq) \land !E(M,\emptyseq)$ ஐ நிறைவுறுத்தும் ஒரு முடிவுறு
  !!{structure}~$\Struct{M''}$ ஐ வரையறுக்கலாம் (பயிற்சி).
  எனினும், அத்தகைய !!{structure}s விலக்கப்படுமாறு~$!T(M,w)$ ஐப்
  பொருத்தமான வகையில் மாற்றலாம்.

  \begin{prob}
    ஒற்றை நிலை $q_0$ ஐயும் ஒற்றைக் கட்டளை
    $\delta(q_0,\TMblank) = \tuple{q,\TMblank,\TMstay}$ ஐயும் கொண்ட
    ஒரு டியூரிங் பொறியாக $M$ இருக்கட்டும்.
    $\Domain{M''} = \{0, 1, 2\}$,
    $\Assign{\prime}{M''}(0) =
    \Assign{\prime}{M'}(1) = 1$ மற்றும்
    $\Assign{\prime}{M''}(2) = 2$ எனவும்,
    $\Assign{<}{M''} = \{\tuple{0,1}, \tuple{1,1}, \tuple{2,2}\}$
    எனவும் வைக்க. $!T(M,\emptyseq)$, $!E(M,\emptyseq)$ ஆகியவை
    உண்மையாகுமாறு $\Assign{\Obj Q_{q_0}}{M''}$,
    $\Assign{\Obj S_{\TMblank}}{M''}$,
    $\Assign{\Obj S_{\TMendtape}}{M''}$ ஆகியவற்றை வரையறுத்து, அவை
    ஏன் உண்மை என்று விளக்குக. குறிப்பு:
    $\delta(q_0, \TMendtape)$ வரையறுக்கப்படவில்லை என்பதைக் கவனிக்க.
    பின்வரும் இரண்டும்
    \begin{align*}
    & \Obj Q_{q_0}(\num{1}, \num{n}) \land \Obj S_{\TMendtape}(\num{0}, \num{n})
      \land 
      \lforall[x][(\num{0} < x \lif \Obj S_\TMblank(x, \num{n}))] \text{\quad எல்லா $n \in \Nat$ க்கும்}\\
    & \lexists[y][(\Obj Q_{q_0}(\num{0}, y) \land \Obj S_{\TMendtape}(\num{0}, y))]
    \end{align*}
    ~$\Struct{M''}$ இல் உண்மையாக இருப்பதை உறுதிசெய்க.
  \end{prob}
\end{explain}

உள்ளீடு $w = \sigma_{i_1}\dots\sigma_{i_k}$ இல் டியூரிங் பொறி~$M$
செயல்படுவதை விவரிக்கும் !!{sentence}s ஐக் கருதுக:
\begin{enumerate}
\item எண்களையும்~$<$ ஐயும் விவரிக்கும் அடிகோள்கள்
(~$!T(M,w)$ இன் வரையறையில் \olref[rep]{sec} இல் உள்ளதைப்போலவே).
\item உள்ளீட்டு அமைவுநிலையை விவரிக்கும் அடிகோள்கள்:
~$!T(M,w)$ இன் வரையறையில் உள்ளதைப்போலவே.
\item ஓர் அமைவுநிலையிலிருந்து அடுத்த அமைவுநிலைக்குச் செல்லும்
  நிலைமாற்றத்தை விவரிக்கும் அடிகோள்கள்:

பின்வருவனவற்றுக்காக, முன்புபோலவே $!A(x, y)$ இருக்கட்டும்; மேலும்
\[
  !B(y) \ident \lforall[x][(x < y \lif \eq/[x][y])].
\]
என இருக்கட்டும்.
\begin{enumerate}
\item \ollabel{rep-right} ஒவ்வொரு கட்டளை $\delta(q_i, \sigma) =
  \tuple{q_j, \sigma', \TMright}$ க்கும், பின்வரும் !!{sentence}:
\begin{align*}
& \lforall[x][\lforall[y][(
   (\Obj Q_{q_i}(x, y) \land \Obj S_{\sigma}(x, y)) \lif {}]] \\
&\qquad   (\Obj Q_{q_j}(x', y') \land \Obj S_{\sigma'}(x, y') \land
!A(x, y) \land !B(y')))
\end{align*}
\item \ollabel{rep-left} ஒவ்வொரு கட்டளை $\delta(q_i, \sigma) =
  \tuple{q_j, \sigma', \TMleft}$ க்கும், பின்வரும் !!{sentence}
\begin{align*}
& \lforall[x][\lforall[y][
    ((\Obj Q_{q_i}(x', y) \land \Obj S_{\sigma}(x', y)) \lif {}]]\\
& \qquad   (\Obj Q_{q_j}(x, y') \land \Obj S_{\sigma'}(x', y') \land
!A(x, y))) \land {}\\
& \lforall[y][((\Obj Q_{q_i}(\num{0}, y) \land \Obj S_{\sigma}(\num{0},
    y)) \lif {}]\\
& \qquad (\Obj Q_{q_j}(\num{0}, y') \land \Obj S_{\sigma'}(\num{0},
  y') \land !A(\num{0}, y) \land !B(y')))
\end{align*}
\item \ollabel{rep-stay} ஒவ்வொரு கட்டளை $\delta(q_i, \sigma) =
  \tuple{q_j, \sigma', \TMstay}$ க்கும், பின்வரும் !!{sentence}:
\begin{align*}
& \lforall[x][\lforall[y][(
   (\Obj Q_{q_i}(x, y) \land \Obj S_{\sigma}(x, y)) \lif {}]] \\
&\qquad (\Obj Q_{q_j}(x, y') \land \Obj S_{\sigma'}(x, y') \land
!A(x, y) \land !B(y')))
\end{align*}
\end{enumerate}
காணக்கூடியவாறு,~$M$ இன் நிலைமாற்றங்களை விவரிக்கும் !!{sentence}s,
$!T(M,w)$ இலுள்ள அவற்றுக்குரிய !!{sentence} ஐப் போலவே உள்ளன;
இறுதியில் $!B(y')$ ஐச் சேர்ப்பது மட்டும் வேறுபாடு. “அடுத்த”
அமைவுநிலையின் எண்~$y'$ ஆனது முந்தைய எல்லா எண்கள்
$\Obj 0$, $\Obj 0'$, \dots ஆகியவற்றிலிருந்தும் வேறுபடுவதை
$!B(y')$ உறுதிசெய்கிறது.
% \item ஒத்திசைவு நிபந்தனைகளை வெளிப்படுத்தும் வாக்கியங்கள்:
% \begin{enumerate}
%   \item\ollabel{rep-con-symb}%
%   எந்தக் கட்டமும் கொடுக்கப்பட்ட எந்த நேரத்திலும் ஒன்றுக்கு மேற்பட்ட
%   குறிகளைத் தன்னில் கொண்டிருக்க முடியாது என்று கூறும் !!^a{sentence}:
%   \[\lforall[x][\lforall[y][(\Obj S_{\sigma}(x, y) \lif \lnot \Obj
%   S_{\sigma'}(x, y))]],\]
%   ~$T$ இன் நாடாக் குறிகளின் ஒவ்வொரு ஜோடி $\sigma \neq \sigma'$ க்கும்.
%   \item\ollabel{rep-con-state}%
%   கொடுக்கப்பட்ட எந்த நேரத்திலும்~$M$ ஒன்றுக்கு மேற்பட்ட நிலைகளில்
%   இருக்க முடியாது என்று கூறும் !!^a{sentence}:
%   \[\lforall[x][\lforall[y][(\Obj Q_{q}(x, y) \lif
%   \lforall[z][\lnot \Obj
%   Q_{q'}(z, y))]],\]
%   ~$T$ இன் நிலைகளின் ஒவ்வொரு ஜோடி $q \neq q'$ க்கும்.
%   \item\ollabel{rep-con-tape}%
%   கொடுக்கப்பட்ட எந்த நேரத்திலும்~$M$ ஒன்றுக்கு மேற்பட்ட நாடாக்
%   கட்டங்களில் இருக்க முடியாது என்று கூறும் !!^a{sentence}:
%   \[\lforall[x][\lforall[y][(\Obj Q_{q}(x, y) \lif
%   \lforall[z][\eq/[z][y]\lif \lnot \Obj
%   Q_{q'}(z, y))]],\]
%   ~$T$ இன் நிலைகளின் ஒவ்வொரு ஜோடி $q$, $q'$ க்கும்.
% \end{enumerate}
\end{enumerate}
டியூரிங் பொறி~$M$ மற்றும் உள்ளீடு~$w$ க்கான மேலுள்ள
!!{sentence}s அனைத்தின் இணையலாக $!T'(M, w)$ இருக்கட்டும்.

\begin{lem}\ollabel{lem:halts-sat}
  உள்ளீடு~$w$ இல் தொடங்கப்பட்ட~$M$ நின்றால்,
  $!T'(M,w) \land !E(M,w)$ கு ஒரு முடிவுறு மாதிரி உள்ளது.
\end{lem}

\begin{proof}
  \olref[ver]{lem:halt-if-valid} இன் நிரூபிப்பிலுள்ளதைப் போலவே
  $\Struct{M'}$ இருக்கட்டும்; பின்வருபவை மட்டும் வேறுபடும்:
  \begin{align*}
    \Domain{M'} & = \{0, \dots, n\},\\
    \Assign{\prime}{M'}(x) & = 
    \begin{cases}
      x + 1 &\text{$x < n$ எனில்}\\
      n &\text{இல்லையெனில்,}
    \end{cases} \\
    \tuple{x,y} \in \Assign{<}{M'} &\text{எனில் மட்டுமே $x < y$ அல்லது $x = y = n$,}
  \end{align*}
  இங்கு $n = \max(k,\len{w})$; மேலும் உள்ளீடு~$w$ இல் தொடங்கப்பட்ட
  ~$M$ ஆனது $k$ படிகளுக்குப் பிறகு நிற்குமாறு அமையும் மிகச்சிறிய
  எண்~$k$. $\Sat{M'}{!T'(M,w) \land E(M,w)}$ என்பதைச் சரிபார்ப்பதைப்
  பயிற்சியாக விட்டுவைக்கிறோம்.
\end{proof}

\begin{prob}
  $\Sat{M'}{!T(M,w) \land E(M,w)}$ என்று நிறுவுவதன் மூலம்,
  \olref[tur][und][tra]{lem:halts-sat} இன் நிரூபிப்பை நிறைவுசெய்க.
\end{prob}

\begin{lem}\ollabel{lem:sat-halts}
  $!T'(M,w) \land !E(M,w)$ கு ஒரு முடிவுறு மாதிரி இருந்தால்,
  உள்ளீடு~$w$ இல் தொடங்கப்பட்ட~$M$ நிற்கிறது.
\end{lem}

\begin{proof}
  எதிர்மாற்றை நிறுவுகிறோம். உள்ளீடு~$w$ இல் தொடங்கப்பட்ட~$M$
  நிற்பதில்லை என்று கொள்க. $!T'(M,w) \land !E(M,w)$ கு எந்த
  மாதிரியுமே இல்லை எனில், முடித்துவிட்டோம். எனவே~$\Struct{M}$
  ஆனது~$!T(M,w) \land !E(M, w)$ இன் ஒரு மாதிரி என்று கருதுக.
  அது முடிவுறுவதாக இருக்க முடியாது என்று காட்ட வேண்டும்.
  
  \olref[ver]{lem:config} இல் உள்ளதைப்போலவே, உள்ளீடு~$w$ இல்
  தொடங்கப்பட்ட~$M$ ஆனது $n$ படிகளுக்குப் பிறகு நிற்கவில்லை எனில்,
  $!T'(M,w) \Entails !C(M, w, n) \land !B(\num{n})$ என்று நிறுவலாம்.
  உள்ளீடு~$w$ இல் தொடங்கப்பட்ட~$M$ நிற்காததால், எல்லா
  ~$n \in \Nat$ க்கும்
  $!T'(M,w) \Entails !C(M, w, n) \land !B(\num{n})$.
  \olref[rep]{prop:mlessk} இன்படி, எல்லா~$k < n$ க்கும்
  $!T'(M,w) \Entails \num{k} < \num{n}$ என்பதைக் கவனிக்க.
  மேலும் $!B(\num{n}) \Entails \num{k} < \num{n} \lif
  \eq/[\num{k}][\num{n}]$. ஆகவே எல்லா~$k < n$ க்கும்
  $\Sat{M}{\eq/[\num{k}][\num{n}]}$; அதாவது, முடிவிலி எண்ணிக்கையிலான
  சொற்கள்~$\num{k}$ அனைத்தும்~$\Struct{M}$ இல் வெவ்வேறு மதிப்புகளைக்
  கொண்டிருக்க வேண்டும். ஆனால் இதற்கு~$\Domain{M}$ முடிவுறாததாக
  இருக்க வேண்டும்; எனவே~$\Struct{M}$ ஆனது
  $!T'(M,w) \land !E(M, w)$ இன் ஒரு முடிவுறு மாதிரியாக இருக்க முடியாது.
\end{proof}

\begin{prob}
  உள்ளீடு~$w$ இல் தொடங்கப்பட்ட~$M$ ஆனது~$n$ படிகளுக்குப் பிறகு
  நிற்கவில்லை எனில், $!T'(M,w) \Entails !B(\num{n})$ என்று
  நிறுவுவதன் மூலம்,
  \olref[tur][und][tra]{lem:sat-halts} இன் நிரூபிப்பை நிறைவுசெய்க.
\end{prob}

\begin{thm}[ட்ராக்டென்ப்ரோட்டின் தேற்றம்]
  \ollabel{thm:trakhtenbrodt}
  முதல்தர தருக்கத்தின் தன்னிச்சையான ஓர் !!{sentence} கு ஒரு
  முடிவுறு மாதிரி உள்ளதா (அதாவது, அது முடிவுறுமாறு நிறைவுறத்தக்கதா)
  என்பது தீர்மானிக்கவியலாதது.
\end{thm}

\begin{proof}
  முடிவுறு நிறைவுறத்தக்கதன்மைச் சிக்கலைத் தீர்மானிக்கும் ஒரு டியூரிங்
  பொறி~$F$ உள்ளது என்று கொள்க. அப்போது எந்த டியூரிங் பொறி~$M$,
  உள்ளீடு~$w$ க்கும் வாக்கியம்~$!T'(M,w) \land !E(M,w)$ ஐக்
  கணித்து, அதற்கு ஒரு முடிவுறு மாதிரி உள்ளதா என்று தீர்மானிக்க~$F$ ஐப்
  பயன்படுத்தலாம். \cref{tur:und:tra:lem:halts-sat,tur:und:tra:lem:sat-halts}
  இன்படி, உள்ளீடு~$w$ இல் தொடங்கப்பட்ட~$M$ நின்றால், எனில் மட்டுமே
  அதற்கு அத்தகைய மாதிரி உள்ளது. ஆகவே தீர்க்கவியலாதது என்று நாம்
  அறிந்த நிற்றல் சிக்கலைத் தீர்க்க~$F$ ஐப் பயன்படுத்தலாம்.
\end{proof}

\begin{cor}\ollabel{cor:fproof-incomp}%
  \emph{முடிவுறு} ஏற்புடைமைக்காகச் சரியானதும் நிறைவானதுமான எந்த
  !!{derivation} முறைமையும் இருக்க முடியாது; அதாவது, ஒவ்வொரு முடிவுறு
  !!{structure}~$\Struct{M}$ க்கும்
  $\Sat{M}{!B}$ எனில் மட்டுமே $\Proves !B$ ஆக உள்ள
  ஓர் !!a{derivation} முறைமை இருக்க முடியாது.
\end{cor}

\begin{proof}
  பயிற்சி.
\end{proof}

\begin{prob}
  \olref[tur][und][tra]{cor:fproof-incomp} ஐ நிறுவுக. ஒவ்வொரு முடிவுறு
  !!{structure} இலும்~$!B$ நிறைவுறுத்தப்படுகிறது எனில் மட்டுமே
  $\lnot !B$ முடிவுறுமாறு நிறைவுறத்தக்கதல்ல என்பதைக் கவனிக்க.
  \olref[tur][und][uns]{thm:valid-ce} இன் பொருளில் முடிவுறு
  நிறைவுறத்தக்கதன்மை ஏன் அரைத் தீர்மானிக்கத்தக்கது என்று விளக்குக.
  முடிவுறு ஏற்புடைமைக்கான ஓர் !!a{derivation} முறைமை இருந்தால்,
  முடிவுறு நிறைவுறத்தக்கதன்மை தீர்மானிக்கத்தக்கதாக இருக்கும் என்று
  வாதிட இதைப் பயன்படுத்துக.
\end{prob}


\end{document}
